\section{A complete existence proof in Lean}
Recall that surjectivity means every target has a preimage. Define it explicitly:
\par\noindent\begin{minipage}{\linewidth}
\begin{lstlisting}
def Surj {A B : Type} (f : A -> B) : Prop :=
  forall y : B, Exists (fun x : A => f x = y)

theorem surj_comp {A B C : Type}
    (f : A -> B) (g : B -> C)
    (hf : Surj f) (hg : Surj g) :
    Surj (fun x => g (f x)) := by
  intro c
  cases hg c with
  | intro b hb =>
    cases hf b with
    | intro a ha =>
      refine Exists.intro a ?_
      show g (f a) = c
      rw [ha, hb]
\end{lstlisting}
\end{minipage}\par

The first line, \code{intro c}, starts with an arbitrary target $c$ in $C$. At this point the goal is an existence statement: find some $a$ in the input type $A$ with $g(f(a))=c$. We do not yet know which $a$ to use. First use the hypothesis about $g$, because $g$ is the map whose outputs lie in $C$ and so can reach our target $c$.

The evidence \code{hg c} says that some input of $g$ reaches $c$; it contains a witness together with a property. The lines
\par\noindent\begin{minipage}{\linewidth}
\begin{lstlisting}
 cases hg c with
 | intro b hb =>
\end{lstlisting}
\end{minipage}\par

mean: unpack that existence evidence, name the witness \code{b}, name its property \code{hb}, and continue inside the indented branch. Here the word \code{intro} after the vertical bar is the name of the way existence evidence is built (a witness together with a property); it is not the standalone tactic \code{intro} that we used to assume a premise. Both uses introduce new names, but they appear in different places in the syntax.

The second unpacking does the same with \code{hf b}. Inside its branch, the available information is
\[
 b\in B,\quad g(b)=c,\qquad a\in A,\quad f(a)=b.
\]
We now have a candidate input, $a$. The line \code{refine Exists.intro a ?\_} proposes $a$ as the witness for the existence goal and leaves the proof of its property, marked by the placeholder \code{?\_}, as the new goal. Lean displays that goal as
\par\noindent\begin{minipage}{\linewidth}
\begin{lstlisting}
(fun x => g (f x)) a = c
\end{lstlisting}
\end{minipage}\par
which still contains the unevaluated application discussed at the end of Chapter~15. The line \code{show g (f a) = c} restates the goal in its beta-reduced form; Lean accepts this because the two forms are the same proposition. This step is not decoration. The next tactic, \code{rw}, rewrites by searching for the literal expression \code{f a}, and that expression only appears once the goal is written in reduced form. Finally, \code{rw [ha, hb]} substitutes $f(a)=b$, turning the goal into $g(b)=c$, and then substitutes $g(b)=c$, turning it into $c=c$, which \code{rw} closes automatically.

The witnesses are chosen in reverse order from function execution: first find an input to $g$ that reaches the final target, then find an input to $f$ that reaches that intermediate value. This is the same argument as in Exercise~\ref{ex:3.2}, now visible in a formal script.

\subsection*{An ordinary proof beside the syntax}
The surjectivity argument can be written without Lean as a three-line dependency chain:
\[
 c\ \longleftarrow\ g(b)=c,\qquad
 b\ \longleftarrow\ f(a)=b,\qquad
 g(f(a))=g(b)=c.
\]
The first two statements use hypotheses to obtain witnesses. The last checks the proposed final witness. In the Lean proof, \code{refine Exists.intro a ?\_} corresponds to proposing the witness, and the last two lines correspond to checking it.

Indentation records which named witnesses are in scope. A witness unpacked inside a branch is available in that branch. The proof does not get to use an unnamed existential object as though its value had already been chosen. This is the formal version of the distinction between ``there is a solution'' and ``let this named object be a solution'' from Chapter~4.
